\section{问题重述及全文大致思路}
围绕低空湍流监测及最优航路规划的需求，三个问题依次解决指标建立、空间融合和未来决策。本文根据题面公开转载中的任务描述整理以下要求\cite{problem}，具体输入字段以原始附件的数据说明为准。
\subsection{问题一}
首先综合利用风廓线雷达与微波辐射计资料，建立不同高度上湍流强度的模型 a。随后限制输入为风廓线雷达资料，从时间、空间或二者结合的角度构建模型 b。以模型 a 为参照优化模型 b，最终输出雷达所在位置的湍流强度垂直廓线。该问题的重点是说明输入减少后哪些信息仍然保留，以及参数标定如何减小两种模型之间的偏差。
\subsection{问题二}
利用地面自动站、风廓线雷达和 S、X 波段多普勒天气雷达构建三维模型 c。输出区域包含全部地面自动站所在矩形范围，高度低于 2 km，水平分辨率为 100 m，垂直分辨率为 50 m，并与提供的地理信息叠加展示。空间网格描述产品的输出尺度，观测分辨率与误差分析用于判断可解释的细节程度。
\subsection{问题三}
利用 02:00 至 05:00 的观测资料及模型 c 形成验证参照，标定基于数值天气预报的模型 d，计算 05:00 至 08:00 的湍流预报；在无数值预报的条件下，利用同一历史观测时段建立非线性外推模型 e，预测 05:00 至 06:00 的湍流场。两条路线均在给定起点、终点和允许飞行区域内，规划仅考虑湍流条件的最优航路。
\subsection{总体解题思路}
要实现低空湍流的空间刻画，需要先将不同观测量放入统一的物理和统计框架。问题一以热力和动力参数为基础，形成可比较的单站湍流指标；问题二将各设备的局地指标转化为带有误差信息的三维场；问题三继续处理时间演化，并把场上的风险转换为图上可累计的路径代价。

对于参数较少、解释关系明确的部分，优先建立线性或物理基线；对于非线性映射和时空外推，再比较树模型、核方法或状态空间模型。所有方法采用相同的时间边界、评价尺度和输出定义。本文的总体解题思路如图\ref{fig:roadmap}所示，各问的输入与交付内容见表\ref{tab:deliverables}。
\samplefigure{roadmap}{全文思路流程图}{fig:roadmap}
\begin{table}[H]\centering\auxtablecaption{三个问题的输入、输出与验证关系}\label{tab:deliverables}
\begin{tabularx}{\linewidth}{lYYY}\toprule
问题 & 输入 & 输出 & 验证参照\\\midrule
问题一 & 雷达与辐射计 & 模型 a、b 及廓线 & a 对 b 的一致性；独立观测检验\\
问题二 & 多设备观测及地图 & 模型 c、三维场和误差场 & 留站或空间块观测\\
问题三 & 历史观测、NWP 和空域 & 模型 d、e 及两类航路 & 历史回报与小图最优路径\\\bottomrule
\end{tabularx}\end{table}
